\begin{abstract}
KI-Coding-Agenten wie Claude Code, Codex oder OpenCode schlagen nicht mehr nur einzelne Codezeilen vor, sondern planen Änderungen selbstständig, führen Werkzeuge aus und bearbeiten ganze Repositories. Damit gewinnt die Gestaltung des Workflows, in den der Agent eingebettet ist, an Bedeutung. Dieses Short Paper untersucht, welche Gestaltungsprinzipien für solche Workflows sich aus der aktuellen Forschung ableiten lassen und welche Hinweise zu ihrem Einfluss auf Codequalität und Produktivität vorliegen. Grundlage ist eine Literaturanalyse von Veröffentlichungen aus dem Zeitraum Januar 2024 bis zum Recherchestichtag am 16.~Oktober 2026. Gesucht wurde mit einem dokumentierten Suchstring und festgelegten Auswahlkriterien in der ACM Digital Library, IEEE Xplore und arXiv. Die Arbeiten stammen aus Software Engineering, Human-Computer Interaction und Machine Learning. Sie werden nach Nutzerstudien, Feldstudien, Benchmark-Evaluationen und Systembeschreibungen kategorisiert und entlang dreier Dimensionen ausgewertet: Autonomiegrad und menschliche Beteiligung, Kontext und Steuerung des Agenten sowie Qualitätssicherung agentisch erzeugten Codes. Die ausgewertete Literatur zeigt unterschiedliche, kontextabhängige Befunde zur Produktivität KI-gestützter Softwareentwicklung, keine durchgängige Verbesserung der Aufgabenerfolgsrate durch Repository-Kontextdateien in den untersuchten Settings sowie Hinweise auf erhöhte Sicherheitsrisiken bei weitreichender Agentenautonomie. Daraus werden praxisorientierte Gestaltungsprinzipien für agentische Entwicklungsworkflows abgeleitet und jeweils mit Literatur begründet. Abschließend werden Forschungslücken diskutiert, insbesondere die begrenzte Aussagekraft gängiger Benchmarks, die begrenzte Langzeitevidenz zur Wartbarkeit in der ausgewerteten Literatur sowie offene Fragen zur Sicherheit und zur veränderten Rolle von Entwicklerinnen und Entwicklern.
\end{abstract}

\begin{IEEEkeywords}
KI-Coding-Agenten, Agentic Coding, Large Language Models, Software Engineering, Entwickler-Workflows, Human-in-the-Loop, Codequalität, Literaturanalyse
\end{IEEEkeywords}
